\chapter{Verificación independiente de evidencia}\label{cap:verificacion}

Cada versión sellada o firmada de un documento puede exportarse como un
paquete de evidencia: un archivo ZIP que reúne el documento y el material
público necesario para comprobarlo. Este capítulo describe su contenido y
las tres vías para verificarlo. La explicación paso a paso, con salidas de
referencia, está en el Anexo~C del reporte técnico (\emph{Guía de
verificación independiente con OpenSSL},
\archivo{latex/chapters/anexo-c-verificacion.tex}); el procedimiento para
quien recibe el paquete está en el Manual de usuario, capítulos
\emph{Validar evidencia} y \emph{Verificación independiente sin cuenta}.

\importante{Un resultado correcto demuestra integridad y correspondencia
técnica dentro de la confianza del prototipo, cuya CA y TSA son internas.
No acredita identidad civil ni sustituye la constancia de un prestador de
servicios de certificación autorizado. Incluir una raíz en el ZIP no la
vuelve confiable: su autenticidad debe confirmarse por un canal
independiente.}

\section{Contenido del paquete}\label{sec:ver-contenido}

El paquete es un ZIP de entradas planas, sin directorios y sin compresión
(método almacenado), con salida determinista y nombres restringidos a un
subconjunto ASCII. El formato se justifica en
\archivo{docs/adr/0006-evidence-package-format.md}. Por ello el validador
del navegador solo admite el archivo tal como se exportó: un ZIP
descomprimido y vuelto a comprimir deja de reconocerse. El paquete se
obtiene con \texttt{despacho-cli package export} (\cref{cap:cli}) o con la
descarga autorizada por expediente
\texttt{GET /api/v1/cases/\{case\_id\}/documents/\{document\_id\}/evidence}
(\cref{cap:api}).

\begin{table}[H]
  \centering
  \caption{Entradas del paquete de sellado del servidor.}
  \label{tab:ver-contenido}
  \begin{tabularx}{\textwidth}{L{4.3cm} Y}
    \toprule
    \textbf{Entrada} & \textbf{Contenido} \\
    \midrule
    \texttt{<documento>} & Documento original, byte a byte. \\
    \texttt{<documento>.sig} & Firma separada RSA PKCS\#1 v1.5 sobre
      SHA-256. \\
    \texttt{<documento>.tsr} & Sello de tiempo RFC~3161 sobre el resumen
      SHA-256 del documento. \\
    \texttt{certificado.pem} & Certificado X.509 del firmante. \\
    \texttt{ca.pem} & Certificado de la autoridad emisora. \\
    \texttt{crl.pem} & Lista de revocación capturada con la evidencia;
      exportar no la renueva. \\
    \texttt{tsa-chain.pem} & Cadena de la autoridad de sellado, cuando
      estaba disponible. En su ausencia el ancla del sello es
      \texttt{ca.pem}. \\
    \texttt{INSTRUCCIONES.md} & Instructivo con los comandos ya
      instanciados para ese paquete. \\
    \bottomrule
  \end{tabularx}
\end{table}

El paquete de una \textbf{firma documental individual} conserva esos
nombres (la firma y el certificado pertenecen al titular) y agrega
\texttt{certificate-binding.bin} y \texttt{.sig} (vínculo entre
certificado y cuenta), \texttt{binding-root.der} y
\texttt{binding-crl.der} (confianza capturada al registrar el vínculo),
\texttt{document-consent.bin} y \texttt{.sig} (consentimiento sobre la
versión exacta) e \texttt{individual-signing.json} (metadatos públicos de
lectura). Cuando la aceptación selló un manifiesto, siguen cuatro entradas:

\begin{itemize}
  \item \texttt{manifiesto.json}: JSON canónico (RFC~8785) que enumera
        cada archivo cubierto con su función, tamaño y huella SHA-256.
  \item \texttt{manifiesto.json.tsr}: sello RFC~3161 sobre la huella del
        manifiesto.
  \item \texttt{tsa-crl.pem}: lista de revocación del emisor del
        certificado de la TSA, capturada al obtener ese sello.
  \item \texttt{recibo-aceptacion.json}: repite la huella del manifiesto
        y el instante de aceptación registrado por la aplicación.
\end{itemize}

Ningún sello cubre el recibo, \texttt{individual-signing.json} ni el
instructivo, y el manifiesto no los enumera. El orden exacto de las
entradas consta en \archivo{docs/http-api.md}; el formato del manifiesto,
en \archivo{docs/adr/0085-in-browser-evidence-validation.md}. Los
contratos bilaterales exportan su propio paquete, con un instructivo
generado desde
\archivo{crates/application/src/contracts/INSTRUCCIONES-CONTRATO.template.md}.

\section{Vías de verificación}\label{sec:ver-vias}

\begin{table}[H]
  \centering
  \caption{Vías para verificar la evidencia y de qué depende cada una.}
  \label{tab:ver-vias}
  \begin{tabularx}{\textwidth}{L{3.1cm} Y Y}
    \toprule
    \textbf{Vía} & \textbf{Qué hace} & \textbf{De qué depende} \\
    \midrule
    En la aplicación
      & El servidor verifica la versión almacenada y reporta integridad,
        firma, certificado y revocación, y sello de tiempo
        (\texttt{POST .../verify}).
      & Sesión con acceso al expediente y el servidor en operación. \\
    Validador del navegador
      & Procesa el ZIP en el equipo del usuario, sin enviarlo al
        servidor, y entrega un resultado global (\emph{Completo},
        \emph{Parcial}, \emph{Fallido} o \emph{No compatible}) y uno por
        componente.
      & Sesión en Qadra y confianza en el programa entregado al
        navegador; las raíces vienen fijas en la versión publicada. \\
    OpenSSL
      & Cuatro comprobaciones con herramientas de terceros.
      & Solo el ZIP, \texttt{unzip} y \texttt{openssl}; sin red ni
        software del proyecto. \\
    \bottomrule
  \end{tabularx}
\end{table}

El validador del navegador ejecuta el código de
\archivo{crates/evidence-verify}, compilado a WebAssembly mediante
\archivo{crates/evidence-verify-wasm} y
\archivo{scripts/build-evidence-verify.sh}. Consulta
\texttt{GET /api/v1/evidence/trust} únicamente para contrastar las raíces
que admite el servidor con las fijas; una raíz que el servidor devuelve
sin estar en esa lista deja el resultado en \emph{Parcial}. Como aplica
perfiles estrictos, su veredicto sobre un paquete histórico puede diferir
del que dio el servidor.

\section{Secuencia mínima con OpenSSL}\label{sec:ver-openssl}

La versión de referencia es OpenSSL~3.5.7; el instructivo de cada paquete
cita la versión de la máquina que lo exportó. Los comandos se ejecutan
dentro del directorio extraído. En los ejemplos el documento se llama
\texttt{document.txt}, como en la ejecución de referencia del Anexo~C;
sustitúyalo por el nombre real.

\begin{verbatim}
unzip -o evidencia.zip -d evidencia
cd evidencia

# 1. Integridad: comparar con el resumen de INSTRUCCIONES.md
openssl dgst -sha256 document.txt

# 2. Firma: debe responder "Verified OK"
openssl x509 -in certificado.pem -pubkey -noout > firmante.pub.pem
openssl dgst -sha256 -verify firmante.pub.pem \
    -signature document.txt.sig document.txt

# 3. Certificado y revocación: "certificado.pem: OK"
cat ca.pem crl.pem > ca-y-crl.pem
openssl verify -crl_check -CAfile ca-y-crl.pem certificado.pem

# 4. Sello de tiempo: "Verification: OK"
openssl ts -verify -data document.txt -in document.txt.tsr \
    -CAfile tsa-chain.pem
\end{verbatim}

Si el paquete no trae \texttt{tsa-chain.pem}, el paso~4 usa
\texttt{-CAfile ca.pem}. Para leer el instante aseverado por la TSA
(\texttt{genTime}) y el resumen que cubre el sello:

\begin{verbatim}
openssl ts -reply -in document.txt.tsr -text
\end{verbatim}

Alterar un byte del documento hace fallar los pasos 1, 2 y 4; el paso~3 no
depende del contenido. La suite fija este comportamiento en
\archivo{crates/bin/tests/package_cli.rs}.

\section{Verificación histórica}\label{sec:ver-historica}

La CRL del paquete es la capturada con la evidencia y tiene fecha de
caducidad. \texttt{openssl verify} evalúa por defecto en la fecha de
ejecución, de modo que el paso~3 puede rechazar una lista que ya superó
su próxima actualización aunque la firma se haya producido con el
certificado vigente. Primero consulte la ventana de la lista:

\begin{verbatim}
openssl crl -in crl.pem -noout -lastupdate -nextupdate
\end{verbatim}

Si la lista ya venció, repita la comprobación en el instante de la firma,
expresado en segundos Unix, con la opción \texttt{-attime}:

\begin{verbatim}
openssl verify -attime <segundos-unix> -crl_check \
    -CAfile ca-y-crl.pem certificado.pem
\end{verbatim}

En el sellado del servidor, el instante de referencia es el
\texttt{genTime} del sello verificado en el paso~4. En la firma
individual, el instructivo del paquete ya trae los comandos con sus dos
instantes (registro del vínculo y aceptación), que son capturas auditadas
de la aplicación y no instantes autenticados por una TSA; si hay
manifiesto, añade la verificación de su sello y de \texttt{tsa-crl.pem}.

\nota{Registre el instante y la política con que evaluó. La verificación
histórica reproduce el estado capturado; el estado actual del certificado
se consulta por separado, con una lista vigente solicitada a la autoridad
emisora (\cref{cap:ca}). La cadena de la TSA se evalúa con la política
vigente de OpenSSL y puede rechazarse si su certificado expiró.}
